\scanpage{11}
\mspage{1}

Hyperrecouvrements.

Soit $\mathcal X$ une catégorie $\mathcal U$-topos. On veut définir \deleted{\unsure{une}} la notion d'\uline{hyperrecouvrement} de $\mathcal X$ (ou encore, du topos défini par $\mathcal X$), indexé par une petite catégorie $A$. \deleted{Un tel hyp} Ce sera par définition un foncteur
\diagram{\begin{tikzcd}[ampersand replacement=\&,row sep=small]
\varphi:\quad A\arrow[r]\&\mathcal X\\
\alpha\arrow[r,mapsto]\&U_\alpha\ (=\varphi(\alpha))
\end{tikzcd}}
, satisfaisant certaines conditions, inspirées par la \unsure{coho} \inserted{``\unsure{Cechiste}''} des recouvrements et de leurs nerfs, et le raffinement de ce \unsure{p.d.v.|p.d.} suggéré par Cartan (pour les espaces topologiques non paracompacts). Voici les conditions qui me viennent à l'esprit:

\uline{Hyp 1} Les objets $U_\alpha$ ($\alpha\in\mathrm{Ob}\,A$) recouvrent l'objet final $e_{\mathcal X}$ de $\mathcal X$.

\deleted{\uline{Hyp 1'} Pour tout $\alpha_0\in\mathrm{Ob}\,A$, le foncteur
\[A/\alpha_0\longrightarrow\mathcal X/U_{\alpha_0}\]
satisfait à Hyp 1, l'ens. des flèches $U_\alpha\longrightarrow U_{\alpha_0}$, associé aux flèches $\alpha\longrightarrow\alpha_0$ de but $\alpha_0$ dans $A$,}
\margin{Est satisfait tautologiquement!}
\margin{10}
